\color{unimain}{\subsection{Логика отключения}}
\color{black}

\begin{enumerate}

\item
Устройство формирует выходное воздействие на  \textit{«Откл. осн. защ.»} по схеме «ИЛИ» в следующих 
случаях: 

\begin{itemize}
    \item Срабатывании ДТЗ;  
    \item Срабатывании ДТЗ НП;
    \item Срабатывании отключающей ступени ЗВП. 
\end{itemize}

\item
Устройство формирует выходное воздействие на  \textit{«Откл. рез. защ.»} по схеме «ИЛИ» в следующих 
случаях: 

\begin{itemize}
    \item Срабатывании МТЗ; 
    \item Срабатывании ТЗНП; 
    \item Срабатывании ТЗОП;
    \item Срабатывании ЗПВГ.
\end{itemize}


\item
Устройство формирует выходное воздействие на \textit{«Откл. реж. защ.»} по схеме «ИЛИ» в 
следующих случаях:

\begin{itemize}
    \item Срабатывание отключающей ступени ЗПВГ; 
    \item Срабатывание ЗМН;
    \item Срабатывание отключающей ступени ЗПН. 
\end{itemize}

\item
Устройство также формирует выходное воздействие на  \textit{«Откл. внеш. защ.»} при появлении сигнала аварийного уровня изоляции ТТ БСК.

\item
Алгоритм работы логики отключения БСК приведен на рисунке \ref{pict:LO}.

\medskip
\begin{figure}[H]
\centering
\includegraphics[width=1\textwidth,height=1\textheight,keepaspectratio]{\fbpath/09.01 БСК/9901 Логика отключения/img/LO.pdf}
\captionof{figure}{Функциональная схема логики отключения БСК}\label{pict_lo:LO}
\end{figure}

\end{enumerate} 